\documentclass{beamer}
\usecolortheme{solarized}
\usepackage{hyperref}
\setbeamertemplate{section in toc}[sections numbered]

\title{Git In Here}
\subtitle{An Unoffical Introduction To Git Version Control Management}
\author{Sam Dixon\\
{\tt \href{mailto:s.dixon@napier.ac.uk}{s.dixon@napier.ac.uk}}}
\date{2018}
\titlegraphic{\includegraphics[height=3.5cm]{../res/git_icon_orange.png}} 

\begin{document}

\frame{\titlepage}

\begin{frame}{Contents}
    \tableofcontents
\end{frame}

\section{What is Version Control?}
\begin{frame}[allowframebreaks]{\secname}
    Version Control Systems are tools that let you track changes

    A number of different VCSs are available:
    \begin{itemize}
        \item Subversion
        \item Mercurial
        \item CVS
        \item Fossil
        \item Git
    \end{itemize}
    Each VCS has its own quirks in style and implementation

\framebreak
    
    Common features of VCSs include:
    \begin{itemize}
        \item Timeline of changes
        \item File reversion  
        \item Branching
        \item Tagging
        \item Collaboration
    \end{itemize}
\end{frame}


\begin{frame}{Why you should use Version control}
    VCSs can help with common problems like:
    \begin{itemize}
        \item Creating multiple copies files ``in case you need the old one''
        \item Put work onto a USB drive to move it between computers
        \item Copied code into email/IM to send it to a friend
        \item Lost work due to loss of power/faulty hard drive
        \item Overwritten someone else's work when collaborating
    \end{itemize}
\end{frame}

\section{What is Git?}
\begin{frame}[allowframebreaks]{\secname}
    Git was developed in 2005 by Linus Torvalds.

    It has become one of the most popular VCSs in use today

    Some companies that use Git include:
    \begin{itemize}
        \item Google
        \item Apple
        \item Microsoft
        \item Facebook
        \item Amazon
        \item Adobe
        \item Mozilla
        \item NASA
    \end{itemize}

\framebreak
    
    Git has a number of key features that make it an attractive VCS
    \begin{itemize}
        \item Distributed
        \item Decentralised
        \item Free (like beer {\it and} like speech)
        \item Multiplatform
        \item Speed
        \item Popularity
    \end{itemize}
\end{frame}

\begin{frame}{Common Terms}
    Git has a number of key terms that you should keep in mind
    \begin{itemize}
        \item Repository 
        \item Commit
        \item Push/Pull
        \item Tag
        \item Branch
        \item Merge
    \end{itemize}
    Some of terms may be used in other VCSs, but their functions may not be
    the same
\end{frame}

\section{Setting up Git}
\begin{frame}[allowframebreaks]{\secname}
    Git can be installed in a number of ways:
    \begin{itemize}
        \item Download directly: {\tt \url{https://git-scm.com/downloads}}
        \item Install a Git visualiser (GitHub Desktop, Sourcetree, Kraken)
        \item Install via a package manager 
    \end{itemize}

\framebreak
    
    Git requires a name and email to create commits \\~\\

    {\tt \$ git congif --global user.name "Sam"} 

    {\tt \$ git config --global user.email sam@mail.com} \\~\\

    These settings are not used as login credentials / tokens
\end{frame}

\section{Git Workflow}
\begin{frame}{\secname}
    The standard process for using Git is as follows:    
    \begin{enumerate}
        \item {\bf Initialise} a Git repo
        \item Do some work (edit files)
        \item {\bf Stage} the changes you want to keep
        \item Finalise the changes in a {\bf Commit}
        \item Repeat from step 2
    \end{enumerate}
\end{frame}

\section{Basic Commands}
\begin{frame}{Getting a Repo}
    Git can create a repo in  a new folder:\\~\\
    {\tt \$ git init project} \\~\\
    Or in an existing directory: \\~\\
    {\tt \$ git init} \\~\\
    It doesn't matter if there are files already in the directory
\end{frame}

\begin{frame}{Staging Files}
    Git will not track files by default - you have to specify: \\~\\
    {\tt \$ git add file\_name} \\~\\
    You can add multiple files in a single command: \\~\\
    {\tt \$ git add file\_1 file\_2} \\~\\
    Or you can use glob wildcards: \\~\\
    {\tt \$ git add *.txt}
\end{frame}

\begin{frame}{Checking File State}
    You can check the state of a repo with the following: \\~\\
    {\tt \$ git status} \\~\\
    Files can be in multiple states:
    \begin{itemize}
        \item Untracked
        \item Modified
        \item Staged
        \item Deleted\\~\\
    \end{itemize}
    You can see how a tracked file has been modified: \\~\\
    {\tt \$ git diff file\_name}
\end{frame}

\begin{frame}{Committing Changes}
    When you are ready to finalise your changes; commit: \\~\\
    {\tt \$ git commit} \\~\\
    You can also declare a commit message in the same command: \\~\\
    {\tt \$ git commit -m "Your message here"}
\end{frame}

\begin{frame}{Viewing History}
    Once you have a few commits, you can review them:\\~\\
    {\tt \$ git log} \\~\\
    View what changes were made in each commit with the {\tt patch} flag: \\~\\
    {\tt \$ git log -p} \\~\\
    There are multiple flags that can alter how the log looks: \\~\\
    {\tt \$ git log --oneline --graph --decorate}
\end{frame}

\begin{frame}{Ignoring Files}
    You may not want to track all the files in repo

    Config files, binary files and build files are often unwanted \\~\\

    You can instruct git to ignore file using a git-ignore file:

    Create a file called .gitignore in the root of your repo

    Add glob patterns for files you want to ignore \\~\\

    {\tt \slash bin} - ignores all files in the bin dir

    {\tt *.o} - ignores all object files \\~\\

    Many sample git-ignore files are available online

\end{frame}

\section{Undoing}
\begin{frame}{\secname}
    The {\tt amend} flag allows you to alter a commit

    It adds your staged files to the {\it previous} commit 

    Use this to add files you forgot, or for quick alterations:\\~\\
    {\tt \$ git commit\\
    \$ git add forgotten\_file\\
    \$ git commit --amend} \\~\\

    When you amend a commit you can also change it's message
\end{frame}

\begin{frame}{Unstaging a File}
    You may accidentally stage a file that you do not want to commit

    You can unstage the file, without losing the changes you made:\\~\\
    {\tt \$ git reset HEAD file\_name} \\~\\
    {\tt HEAD} is Git's name for the tip of the current branch   
\end{frame}

\begin{frame}{Unmodifying a File}
    You may edit a file and decide to not keep the modifications

    You revert a file to how it was when it was last committed: \\~\\
    
    {\tt \$ git checkout -- file\_name} \\~\\

    This command is dangerous - you cant recover uncommitted changes
\end{frame}

\begin{frame}{Reverting a Commit}
    A main feature of a VCS is the ability to undo previous changes \\~\\

    Git provides this function with the {\tt revert} command \\~\\

    First use the {\tt log} command to display commit hashes \\~\\

    You can then undo the changes that a commit introduced: \\~\\
    
    {\tt \$ git revert commit\_hash} \\~\\

    Git will then create a commit inverse of the selected hash \\~\\
\end{frame}

\begin{frame}{Reverting a Series of Commits}
    You can also undo a series  of commits in a single command \\~\\

    Use {\tt log} to find the first and last commit to remove \\~\\

    {\tt \$ git revert from\_hash..to\_hash} \\~\\

    Use the {\tt --no-edit} flag for git if you don't want to change the default
    revert message
\end{frame}

\section{Branching}
\begin{frame}{\secname}
    Branches are a key feature of Git \\~\\

    They allow different versions of a project to exists simultaneously \\~\\
    
    Allows for parallel development of features \\~\\

    Branches can be used to experiment or write patches \\~\\

    Branches are inexpensive to create/delete \\~\\
\end{frame}

\begin{frame}{Creating Branches}
    You can quickly create branches:

    {\tt \$ git branch branch\_name}

    This will create a new branch starting at your current commit\\~\\

    You can list the different branches in a repo:

    {\$ git branch}

    The {\tt status} command will inform you of which branch you are on
\end{frame}

\begin{frame}{Swapping branches}
    Once a branch has been created, you need to switch to it: \\~\\

    {\tt \$ git checkout branch\_name} \\~\\

    Git may stop you from swapping if you have uncommitted changes
\end{frame}

\begin{frame}{Merging Branches}
    You can add the changes from one branch to another

    Merging will make Git apply the divergent commits from a branch \\~\\

    First, swap to the branch you want to apply the commits to:

    {\tt \$ git checkout master}

    Then merge in your branch with the commits you want:

    {\tt \$ git merge branch\_name}
\end{frame}

\begin{frame}{Conflicts}
    Sometimes you will not be able to merge due to conflicts \\~\\

    Conflicts occur when two branches edit a file differently \\~\\

    Conflicts must be resolved before a merge completes \\~\\

    The {\tt status} command informs you which files have problems
\end{frame}

\begin{frame}{Resolving Conflicts}
    A simple way to fix a conflict is to simply open the file \\~\\

    Git will automatically tag the conflicting area \\~\\

    You can edit the file to keep the version you want \\~\\

    When you a happy with the file, stage it and {\tt commit}
\end{frame}

\section{Remotes}
\begin{frame}[allowframebreaks]{\secname}
    Git can be used entirely locally\\~\\

    But remotes allow you to collaborate and store projects\\~\\

    A remote is copy of a Git repo ``somewhere else''\\~\\

    A number of services will host repos for you\\~\\

    You can even run your own Git server\\~\\

\framebreak

    There a two ways to set up repo remotes \\~\\

    You can make a local copy of a remote repo: \\~\\

    {\tt \$ git clone remote\_url} \\~\\

    This will copy the  repo and set the url as the default remote
    \\~\\

    You can also add a remote to a local repo manually: \\~\\

    {\tt \$ git remote add remote\_name remote\_url} \\~\\

    The default name for a repo's primary remote is {\tt origin} - but you can
    make it what ever you want

\framebreak
    
    You can see what remotes are configured to a repo: \\~\\

    {\tt \$ git remote} \\~\\

    To see the url's of theses remotes include the {\tt verbose} flag: \\~\\

    {\tt \$ git remote -v} \\~\\

    You can rename remote:\\~\\
    {\tt \$ git remote rename remote\_name new\_name} \\~\\

    Or remove it if you dont need to use it any more: \\~\\

    {\tt \$ git remote remove remote\_name}

\framebreak

    You can share local commits to a remote: \\~\\

    {\tt \$ git push remote\_name branch\_name} \\~\\

    The default branch name is {\tt master} \\~\\

    You can query a remote to learn its state: \\~\\

    {\tt \$ git remote show remote\_name} \\~\\

\framebreak

    To get commits from a remote repo, you must fetch it:\\~\\

    {\tt \$ git fetch remote\_name} \\~\\

    Fetching will copy data to a remote branch - without changing your
    files\\~\\
    
    You can inspect these branches with commands like {\tt checkout}
    and {\tt log}
    
\framebreak

    You can add changes from a remote branch with merge: \\~\\

    {\tt \$ git merge remote\_branch\_name} \\~\\

    You can revert theses merges if you don't like them
    \end{frame}

\section{Useful Links}
\begin{frame}{\secname}
     Pro Git Book:
    \url{https://git-scm.com/book/en/v2} \\~\\   
    
    Atlassian Git Tutorials:
    \url{https://www.atlassian.com/git/tutorials} \\~\\

    GitHub interactive Tutorial:
    \url{https://try.github.io/levels/1/challenges/1}
\end{frame}

\end{document}

